% GENERATED FILE. Edit canonical lessons, then run npm run generate:books.
% canonical-source-hash: fnv1a64:2dcf27b88c687416
% canonical-lessons: MW-C202-murno, MW-C202-pachho, MW-C202-ulto, MW-C202-ado, MW-C202-tirchho

\chapter{Turn and Go Back}
\label{ch:mw-turn-and-go-back}
\hlchaptermodality{202}

\begin{chapteropening}
\textbf{By the end of this chapter:} \emph{I can follow a turn and find my way back.}

Complete the last lesson of chapter 202: I can follow a turn and find my way back.
\end{chapteropening}

\section[\texorpdfstring{\mw{मुड़णो} (muṛṇo)}{मुड़णो (muṛṇo)}]{\mw{मुड़णो} (muṛṇo) — to turn}

\label{lesson:MW-C202-murno}



\begin{quote}
Before the new one: say the Marwadi for an inn, then the Marwadi for to get in.
\end{quote}

\subsection*{You'll want to know: \mw{मुड़णो}}
\textbf{\mw{मुड़णो}} — \emph{muṛṇo} — \textquotedblleft{}to turn\textquotedblright{}.

\begin{quote}
\mw{आप} \mw{मोड़} \mw{माथै} \mw{मुड़ो।}
\end{quote}

\begin{quote}
\emph{āp moṛ māthai muṛo.}
\end{quote}

\begin{quote}
Turn at the bend.
\end{quote}

Past (he …): \mw{वो} \mw{मुड़्यो}, \emph{vo muṛyo}.

\begin{cousinweb}
\textbf{\mw{मुड़णो}} (\emph{muṛṇo}) is inherited from Sanskrit \emph{muṭati}, \textquotedblleft{}twists\textquotedblright{}. CDIAL gives no Marwari form in this entry, but lists Hindi \emph{muṛnā}, \textquotedblleft{}to be turned\textquotedblright{} (CDIAL 10186). \textbf{\mw{मोड़}}, \textquotedblleft{}a bend\textquotedblright{}, belongs to the same family.
\end{cousinweb}

\begin{grammarlens}[title={What you've built}]
One more word for this chapter.
\end{grammarlens}

\subsection*{Your turn}
\begin{itemize}
\raggedright
  \item \emph{Say it:} \emph{muṛṇo}
  \item \emph{Say it:} \emph{muṛṇo}, once more
  \item \emph{Say it:} \emph{āp moṛ māthai muṛo.} — Turn at the bend.
  \item \emph{Recall:} say the Marwadi for an inn, then the Marwadi for to get in, then say \emph{muṛṇo} again
\end{itemize}

\begin{tcolorbox}[breakable,colback=teal!4,colframe=teal!35!black,title={Before you move on}]
What does \mw{मुड़णो} mean? (\textquotedblleft{}To turn\textquotedblright{}.) Say it once more.
\end{tcolorbox}
\section[\texorpdfstring{\mw{पाछो} (pāchho)}{पाछो (pāchho)}]{\mw{पाछो} (pāchho) — back again; back the way one came}

\label{lesson:MW-C202-pachho}



\begin{quote}
Before the new one: say the Marwadi for to get in, then the Marwadi for to turn.
\end{quote}

\subsection*{You'll want to know: \mw{पाछो}}
\textbf{\mw{पाछो}} — \emph{pāchho} — \textquotedblleft{}back again; back the way one came\textquotedblright{}.

\begin{quote}
\mw{हूं} \mw{काल} \mw{पाछो} \mw{आसूं।}
\end{quote}

\begin{quote}
\emph{hū̃ kāl pāchho āsū̃.}
\end{quote}

\begin{quote}
I'll come back tomorrow.
\end{quote}

\begin{cousinweb}
\textbf{\mw{पाछो}} (\emph{pāchho}) is inherited from Sanskrit \emph{paścāt}, \textquotedblleft{}from behind\textquotedblright{}; CDIAL lists Old Marwari \emph{pācha}, \textquotedblleft{}behind\textquotedblright{} (CDIAL 7999).
\end{cousinweb}

\begin{grammarlens}[title={What you've built}]
One more word for this chapter.
\end{grammarlens}

\subsection*{Your turn}
\begin{itemize}
\raggedright
  \item \emph{Say it:} \emph{pāchho}
  \item \emph{Say it:} \emph{pāchho}, once more
  \item \emph{Say it:} \emph{hū̃ kāl pāchho āsū̃.} — I'll come back tomorrow.
  \item \emph{Recall:} say the Marwadi for to get in, then the Marwadi for to turn, then say \emph{pāchho} again
\end{itemize}

\begin{tcolorbox}[breakable,colback=teal!4,colframe=teal!35!black,title={Before you move on}]
What does \mw{पाछो} mean? (\textquotedblleft{}Back again; back the way one came\textquotedblright{}.) Say it once more.
\end{tcolorbox}
\section[\texorpdfstring{\mw{उलटो} (ulṭo)}{उलटो (ulṭo)}]{\mw{उलटो} (ulṭo) — reversed; the wrong way round}

\label{lesson:MW-C202-ulto}



\begin{quote}
Before the new one: say the Marwadi for to turn, then the Marwadi for back again.
\end{quote}

\subsection*{You'll want to know: \mw{उलटो}}
\textbf{\mw{उलटो}} — \emph{ulṭo} — \textquotedblleft{}reversed; the wrong way round\textquotedblright{}.

\begin{quote}
\mw{म्हारो} \mw{कुरतो} \mw{उलटो} \mw{है।}
\end{quote}

\begin{quote}
\emph{mhāro kurto ulṭo hai.}
\end{quote}

\begin{quote}
My kurta is inside out.
\end{quote}

\begin{cousinweb}
\textbf{\mw{उलटो}} (\emph{ulṭo}) is inherited from a reconstructed \emph{*ullaṭati}, \textquotedblleft{}turns over\textquotedblright{}; CDIAL lists Old Marwari \emph{ulaṭo}, \textquotedblleft{}reversed\textquotedblright{} (CDIAL 2368).
\end{cousinweb}

\begin{grammarlens}[title={What you've built}]
One more word for this chapter.
\end{grammarlens}

\subsection*{Your turn}
\begin{itemize}
\raggedright
  \item \emph{Say it:} \emph{ulṭo}
  \item \emph{Say it:} \emph{ulṭo}, once more
  \item \emph{Say it:} \emph{mhāro kurto ulṭo hai.} — My kurta is inside out.
  \item \emph{Recall:} say the Marwadi for to turn, then the Marwadi for back again, then say \emph{ulṭo} again
\end{itemize}

\begin{tcolorbox}[breakable,colback=teal!4,colframe=teal!35!black,title={Before you move on}]
What does \mw{उलटो} mean? (\textquotedblleft{}Reversed; the wrong way round\textquotedblright{}.) Say it once more.
\end{tcolorbox}
\section[\texorpdfstring{\mw{आडो} (āḍo)}{आडो (āḍo)}]{\mw{आडो} (āḍo) — across; crosswise}

\label{lesson:MW-C202-ado}



\begin{quote}
Before the new one: say the Marwadi for back again, then the Marwadi for the wrong way round.
\end{quote}

\subsection*{You'll want to know: \mw{आडो}}
\textbf{\mw{आडो}} — \emph{āḍo} — \textquotedblleft{}across; crosswise\textquotedblright{}.

\begin{quote}
\mw{रूंख} \mw{सड़क} \mw{माथै} \mw{आडो} \mw{है।}
\end{quote}

\begin{quote}
\emph{rū̃kh saṛak māthai āḍo hai.}
\end{quote}

\begin{quote}
A tree lies across the road.
\end{quote}

\begin{cousinweb}
\textbf{\mw{आडो}} (\emph{āḍo}) is inherited from a reconstructed \emph{*aḍḍa}, \textquotedblleft{}transverse\textquotedblright{}; CDIAL lists Old Marwari \emph{āḍo}, \textquotedblleft{}transverse\textquotedblright{} (CDIAL 189).
\end{cousinweb}

\begin{grammarlens}[title={What you've built}]
One more word for this chapter.
\end{grammarlens}

\subsection*{Your turn}
\begin{itemize}
\raggedright
  \item \emph{Say it:} \emph{āḍo}
  \item \emph{Say it:} \emph{āḍo}, once more
  \item \emph{Say it:} \emph{rū̃kh saṛak māthai āḍo hai.} — A tree lies across the road.
  \item \emph{Recall:} say the Marwadi for back again, then the Marwadi for the wrong way round, then say \emph{āḍo} again
\end{itemize}

\begin{tcolorbox}[breakable,colback=teal!4,colframe=teal!35!black,title={Before you move on}]
What does \mw{आडो} mean? (\textquotedblleft{}Across; crosswise\textquotedblright{}.) Say it once more.
\end{tcolorbox}
\section[\texorpdfstring{\mw{तिरछो} (tirchho)}{तिरछो (tirchho)}]{\mw{तिरछो} (tirchho) — slanting}

\label{lesson:MW-C202-tirchho}



\begin{quote}
Before the new one: say the Marwadi for the wrong way round, then the Marwadi for crosswise.
\end{quote}

\subsection*{You'll want to know: \mw{तिरछो}}
\textbf{\mw{तिरछो}} — \emph{tirchho} — \textquotedblleft{}slanting\textquotedblright{}.

\begin{quote}
\mw{ये} \mw{गेलो} \mw{तिरछो} \mw{है।}
\end{quote}

\begin{quote}
\emph{ye gelo tirchho hai.}
\end{quote}

\begin{quote}
This path runs at a slant.
\end{quote}

\begin{cousinweb}
\textbf{\mw{तिरछो}} (\emph{tirchho}) is inherited from Sanskrit \emph{tiraśca}, \textquotedblleft{}slanting, crosswise\textquotedblright{}. CDIAL gives no Marwari form in this entry, but lists Hindi \emph{tirchā}, \textquotedblleft{}slanting\textquotedblright{} (CDIAL 5822).
\end{cousinweb}

\begin{grammarlens}[title={What you've built}]
One more word for this chapter.
\end{grammarlens}

\subsection*{Your turn}
\begin{itemize}
\raggedright
  \item \emph{Say it:} \emph{tirchho}
  \item \emph{Say it:} \emph{tirchho}, once more
  \item \emph{Say it:} \emph{ye gelo tirchho hai.} — This path runs at a slant.
  \item \emph{Recall:} say the Marwadi for the wrong way round, then the Marwadi for crosswise, then say \emph{tirchho} again
\end{itemize}

\begin{tcolorbox}[breakable,colback=teal!4,colframe=teal!35!black,title={Before you move on}]
What does \mw{तिरछो} mean? (\textquotedblleft{}Slanting\textquotedblright{}.) Say it once more.
\end{tcolorbox}
